Habsora ("The Gospel") is an Israeli Defense Forces (IDF) AI-targeting system developed primarily by Unit 8200 within the Military Intelligence Directorate. First deployed during the May 2021 Operation "Guardian of the Walls", it generates target recommendations using a multimodal intelligence pipeline that encompasses signal intelligence, telephone metadata, satellite and drone imagery, social graph analysis, and life-pattern modelling \cite{abraham_2024, dwoskin_2024}. The system is not a single model, but an ensemble of multiple specialised algorithms fused together \cite{dwoskin_2024}. During the 2023-2024 Gaza conflict, Habsora was reported to generate over one hundred target recommendations per day, which in prior years required approximately twelve months of human analysts \cite{abraham_2024, ynet_2024}.

Habsora's AI capability is the probabilistic classification and prioritisation of locations and individuals as military objectives. It assigns confidence scores derived from large-scale pattern matching across a shared intelligence pool, rather than positive identification in the traditional intelligence analysis \cite{ootlaws}. Military commands nominally authorise each recommendation, yet their review window documented for each target was approximately twenty seconds \cite{abraham_2024}.

This creates a stakeholder dynamic with both ethical and legal significance. It's output influencing both IDF commands, who operate under high cognitive load and time pressure; Palestinian civilians in the target environment bearing the lethal consequences and the broader international community which depends on International Humanitarian Law (IHL) compliance for the legitimacy of armed conflict. \citet{mhcof} offer a framework for this tension. They propose a two-axis model of Meaningful Human Control (MHC) which distinguishes between high automation and genuine human control, which both require adequate information and real freedom of action. The central argument of this report is that Habsora structurally undermines the second condition, and that this gap carries cascading implications across the technical, sociotechnical, and governance layers identified by \citet{info12090385} in their Comprehensive Human Oversight Framework (CHOF).